\documentclass[a4paper,12pt]{report}
\usepackage{usual}
% \usepackage[backend=biber,style=phys,sorting=none]{biblatex}
% \addbibresource{ass2ref.bib}

\usepackage[a4paper, margin=1in]{geometry}

\setlength{\parindent}{0pt}
\setstretch{1.2}

\newcommand{\ti}[1]{\textit{#1}}

% ===============================================================================================================

\title{Notes on Radio Evolution of Galaxies}
\author{Anurag Tamuli \\ PRL '26}
\date{}

\begin{document}

\maketitle

\begin{abstract}
  The primary source for these notes has been \ti{Extragalactic Astronomy and Cosmology} by \ti{Peter Schneider}. These notes were made during my summer internship at PRL Thaltej from May 11 to July 10, 2026.
\end{abstract}

\tableofcontents

\newpage

\chapter{The World of Galaxies}

Ever since the Hubble Deep Field image, the whole of humanity has set their sight to discover what lies far beyond they could ever comprehend. When dealing with extragalactic astronomy, the most notable entities are arguably other galaxies, and this is a chapter all about them.

\begin{figure}[h]
  \centering
  \includegraphics[width=0.8\textwidth]{hdf.jpg}
  \caption{\small The Hubble Deep Field image produced by \ti{Edwin Hubble} changed the world forever. All the dots in the image are galaxies, not stars.}
\end{figure}

\clearpage

\section{Classification of Galaxies}

\subsection{Morphological Classification}
On the basis of shape or morphology, galaxies have been historically classified into the following categories:

\begin{enumerate}
  \item \textbf{Elliptical Galaxies}: Their isophotes\footnote{sdsds} resemble ellipses. They are further subclassified according to the \textbf{ellipticity} of the isophotes: $\epsilon = 1 - b/a$, where a, b are the semi-major and semi-minor axes. Then we define $N = 10 \epsilon$ which gives us the representation EN for elliptical galaxies. An E7 galaxy is an elliptical galaxy with $b/a = 0.3$.
  \item \textbf{Spiral Galaxies}: Their isophotes look like spirals. 
\end{enumerate}

% \printbibliography

\end{document}
